{{دیوان
| عنوان = قید خانے میں معتمدؔ کی فریاد
| ماخذ = ویکی ماخذ
| ماخذ_لنک = https://ur.wikisource.org/wiki/%D9%82%DB%8C%D8%AF_%D8%AE%D8%A7%D9%86%DB%92_%D9%85%DB%8C%DA%BA_%D9%85%D8%B9%D8%AA%D9%85%D8%AF%D8%94%DA%A9%DB%8C_%D9%81%D8%B1%DB%8C%D8%A7%D8%AF
| مقابلہ = اقبال سائبر لائبریری (اقبال اکادمی پاکستان)
| مقابلہ_لنک = https://www.iqbalcyberlibrary.net
}}
معتمدؔ اشبیلیہ کا بادشاہ اور عربی شاعر تھا۔ ہسپانیہ کے ایک حکمران نے اس کو شکست دے کر قید میں ڈال دیا تھا۔ معتمدؔ کی نظمیں انگریزی میں ترجمہ ہوکر “وِزڈم آف دی ایسٹ سِیریز” میں شائع ہو چکی ہیں

<poem>
اک فغانِ بے شرر سینے میں باقی رہ گئی
سوز بھی رُخصت ہُوا، جاتی رہی تاثیر بھی

مردِ حُر زنداں میں ہے بے نیزہ و شمشیر آج
مَیں پشیماں ہوں، پشیماں ہے مری تدبیر بھی

خود بخود زنجیر کی جانب کھِنچا جاتا ہے دل
تھی اسی فولاد سے شاید مری شمشیر بھی

جو مری تیغِ دو دم تھی، اب مری زنجیر ہے
شوخ و بے پروا ہے کتنا خالقِ تقدیر بھی!
</poem>
